\documentclass[11pt]{letter}

\usepackage[utf8]{inputenc}
\usepackage[T1]{fontenc}
\usepackage{geometry}
\usepackage{hyperref}

\geometry{left=2.5cm, right=2.5cm, top=2.5cm, bottom=2.5cm}

\signature{Xiaohui Li}
\address{Independent Researcher\\Baoji, China\\Email: danceluck@163.com\\ORCID: 0009-0003-8233-7191}

\begin{document}

\begin{letter}{
    The Editors\\
    \textit{Nature}\\
    Springer Nature\\
    London, UK
}

\opening{Dear Editors,}

We wish to submit our manuscript entitled \textbf{``Information-Processing Modes Determine the Architecture and Outcome of Aggressive Wars''} for consideration as an Article in \textit{Nature}.

Why do some states initiate aggressive wars while others resolve conflicts peacefully? And why do aggressors so often lose, even when they possess material superiority? Our paper addresses these questions through a novel information-theoretic lens. We introduce the \textit{Information-Granularity Theory} (IGT), which posits that a society's problem-solving capacity is constrained by the smallest information-processing unit (IPU) it can mobilize. When core concepts such as nuclear deterrence or armed coercion become indivisible ``smallest units,'' the system cannot decompose complex problems and instead accumulates unresolved complexity---a process we term the ``coarse-grained trap.'' Eventually, the system either collapses or exports its entropy through aggressive war.

The paper makes three principal contributions:

\begin{enumerate}
    \item \textbf{Theoretical.} We formalize the relationship between information granularity, problem decomposition, and systemic collapse through two theorems: the Decomposability Theorem and the Entropy-Accumulation Collapse Theorem.
    \item \textbf{Empirical.} We construct an Information-Granularity Index (IGI) from V-Dem democracy and press-freedom variables and benchmark it against the Correlates of War (COW) dataset (1816--2007). States with lower IGI initiate aggressive wars more frequently and are significantly more likely to be defeated ($\beta_1=-2.34$, $p<0.01$).
    \item \textbf{Computational.} We develop an agent-based model, calibrate it to seven historical wars, and conduct extensive sensitivity analyses. Across all parameter combinations, fine-grained states show markedly lower aggression rates and higher wartime win rates.
\end{enumerate}

The work bridges information theory, computational social science, and international relations. It offers a fresh explanation for the democratic peace and generates concrete, testable predictions about state behavior. All data and code are openly available at \url{https://github.com/xiaohui-li-lab/}.

The manuscript is original, has not been published elsewhere, and is not under consideration by another journal. All authors have approved the final version and agree to its submission.

We thank you for your consideration and look forward to your response.

\closing{Sincerely,}

\vspace{1cm}
Xiaohui Li\\
Independent Researcher\\
ORCID: 0009-0003-8233-7191

\end{letter}
\end{document}
